% GENERATED FILE. Edit canonical lessons, then run npm run generate:books.
% canonical-source-hash: fnv1a64:56611eaeb0427011
% canonical-lessons: JA-C114-hakusai, JA-C114-ninniku, JA-C114-anko, JA-C114-hachimitsu, JA-C114-ami

\chapter{Chinese cabbage, Garlic, Sweet bean paste, Honey, Net}
\label{ch:ja-chinese-cabbage-garlic-sweet-bean-paste-honey-net}
\hlchaptermodality{114}

\begin{chapteropening}
\textbf{By the end of this chapter:} \emph{I can say Chinese cabbage, garlic, sweet bean paste, honey and a net.}

Complete the last lesson of chapter 114: I can say Chinese cabbage, garlic, sweet bean paste, honey and a net.
\end{chapteropening}

\section[\texorpdfstring{\ja{はくさい} (hakusai)}{はくさい (hakusai)}]{\ja{はくさい} (hakusai) — Chinese cabbage}

\label{lesson:JA-C114-hakusai}



\begin{quote}
Before the new one: say the Japanese for a bonito, then the Japanese for wakame seaweed.
\end{quote}

\subsection*{You'll want to know: \ja{はくさい}}
\textbf{\ja{はくさい}} — \emph{hakusai} — \textquotedblleft{}Chinese cabbage\textquotedblright{}.

\begin{grammarlens}[title={What you've built}]
One more word for this chapter.
\end{grammarlens}

\subsection*{Your turn}
\begin{itemize}
\raggedright
  \item \emph{Say it:} \emph{hakusai}
  \item \emph{Say it:} \emph{hakusai}, once more
  \item \emph{Say it:} say \emph{hakusai}
  \item \emph{Recall:} say the Japanese for a bonito, then the Japanese for wakame seaweed, then say \emph{hakusai} again
\end{itemize}

\begin{tcolorbox}[breakable,colback=teal!4,colframe=teal!35!black,title={Before you move on}]
What does \ja{はくさい} mean? (\textquotedblleft{}Chinese cabbage\textquotedblright{}.) Say it once more.
\end{tcolorbox}
\section[\texorpdfstring{\ja{にんにく} (ninniku)}{にんにく (ninniku)}]{\ja{にんにく} (ninniku) — garlic}

\label{lesson:JA-C114-ninniku}



\begin{quote}
Before the new one: say the Japanese for wakame seaweed, then the Japanese for Chinese cabbage.
\end{quote}

\subsection*{You'll want to know: \ja{にんにく}}
\textbf{\ja{にんにく}} — \emph{ninniku} — \textquotedblleft{}garlic\textquotedblright{}.

\begin{grammarlens}[title={What you've built}]
One more word for this chapter.
\end{grammarlens}

\subsection*{Your turn}
\begin{itemize}
\raggedright
  \item \emph{Say it:} \emph{ninniku}
  \item \emph{Say it:} \emph{ninniku}, once more
  \item \emph{Say it:} say \emph{ninniku}
  \item \emph{Recall:} say the Japanese for wakame seaweed, then the Japanese for Chinese cabbage, then say \emph{ninniku} again
\end{itemize}

\begin{tcolorbox}[breakable,colback=teal!4,colframe=teal!35!black,title={Before you move on}]
What does \ja{にんにく} mean? (\textquotedblleft{}Garlic\textquotedblright{}.) Say it once more.
\end{tcolorbox}
\section[\texorpdfstring{\ja{あんこ} (anko)}{あんこ (anko)}]{\ja{あんこ} (anko) — sweet bean paste}

\label{lesson:JA-C114-anko}



\begin{quote}
Before the new one: say the Japanese for Chinese cabbage, then the Japanese for garlic.
\end{quote}

\subsection*{You'll want to know: \ja{あんこ}}
\textbf{\ja{あんこ}} — \emph{anko} — \textquotedblleft{}sweet bean paste\textquotedblright{}.

\begin{grammarlens}[title={What you've built}]
One more word for this chapter.
\end{grammarlens}

\subsection*{Your turn}
\begin{itemize}
\raggedright
  \item \emph{Say it:} \emph{anko}
  \item \emph{Say it:} \emph{anko}, once more
  \item \emph{Say it:} say \emph{anko}
  \item \emph{Recall:} say the Japanese for Chinese cabbage, then the Japanese for garlic, then say \emph{anko} again
\end{itemize}

\begin{tcolorbox}[breakable,colback=teal!4,colframe=teal!35!black,title={Before you move on}]
What does \ja{あんこ} mean? (\textquotedblleft{}Sweet bean paste\textquotedblright{}.) Say it once more.
\end{tcolorbox}
\section[\texorpdfstring{\ja{はちみつ} (hachimitsu)}{はちみつ (hachimitsu)}]{\ja{はちみつ} (hachimitsu) — honey}

\label{lesson:JA-C114-hachimitsu}



\begin{quote}
Before the new one: say the Japanese for garlic, then the Japanese for sweet bean paste.
\end{quote}

\subsection*{You'll want to know: \ja{はちみつ}}
\textbf{\ja{はちみつ}} — \emph{hachimitsu} — \textquotedblleft{}honey\textquotedblright{}.

\begin{grammarlens}[title={What you've built}]
One more word for this chapter.
\end{grammarlens}

\subsection*{Your turn}
\begin{itemize}
\raggedright
  \item \emph{Say it:} \emph{hachimitsu}
  \item \emph{Say it:} \emph{hachimitsu}, once more
  \item \emph{Say it:} say \emph{hachimitsu}
  \item \emph{Recall:} say the Japanese for garlic, then the Japanese for sweet bean paste, then say \emph{hachimitsu} again
\end{itemize}

\begin{tcolorbox}[breakable,colback=teal!4,colframe=teal!35!black,title={Before you move on}]
What does \ja{はちみつ} mean? (\textquotedblleft{}Honey\textquotedblright{}.) Say it once more.
\end{tcolorbox}
\section[\texorpdfstring{\ja{あみ} (ami)}{あみ (ami)}]{\ja{あみ} (ami) — a net}

\label{lesson:JA-C114-ami}



\begin{quote}
Before the new one: say the Japanese for sweet bean paste, then the Japanese for honey.
\end{quote}

\subsection*{You'll want to know: \ja{あみ}}
\textbf{\ja{あみ}} — \emph{ami} — \textquotedblleft{}a net\textquotedblright{}.

\begin{grammarlens}[title={What you've built}]
One more word for this chapter.
\end{grammarlens}

\subsection*{Your turn}
\begin{itemize}
\raggedright
  \item \emph{Say it:} \emph{ami}
  \item \emph{Say it:} \emph{ami}, once more
  \item \emph{Say it:} say \emph{ami}
  \item \emph{Recall:} say the Japanese for sweet bean paste, then the Japanese for honey, then say \emph{ami} again
\end{itemize}

\begin{tcolorbox}[breakable,colback=teal!4,colframe=teal!35!black,title={Before you move on}]
What does \ja{あみ} mean? (\textquotedblleft{}A net\textquotedblright{}.) Say it once more.
\end{tcolorbox}
